% Generated draft: requires mathematical and source review.
\documentclass[11pt]{article}
\usepackage[margin=25mm]{geometry}
\usepackage{fontspec}
\setmainfont{Latin Modern Roman}
\usepackage{amsmath,amssymb,mathtools,unicode-math}
\setmathfont{Latin Modern Math}
\usepackage{enumitem,needspace,xurl,hyperref}
\setlength{\parindent}{0pt}
\setlength{\parskip}{5pt}
\setlength{\emergencystretch}{3em}
\newcommand{\N}{\mathbb{N}}
\newcommand{\Z}{\mathbb{Z}}
\newcommand{\Q}{\mathbb{Q}}
\newcommand{\R}{\mathbb{R}}
\newcommand{\C}{\mathbb{C}}
\newcommand{\F}{\mathbb{F}}
\newcommand{\A}{\mathbb{A}}
\newcommand{\PP}{\mathbb{P}}
\newcommand{\E}{\mathbb{E}}
\newcommand{\eps}{\varepsilon}
\newcommand{\dd}{\,\mathrm{d}}
\begin{document}
\section*{UnsolvedMath ID: 700531}
\textbf{Source problem code:} catalogue entry
\subsection{Definitions and mathematical statement}
The supplied statement depends on the notation of Kourovka Notebook Problem 8.82, which is not included in the record. In particular, the object \(\pi\), the quantity \(r(\pi)\), and their domains are not defined here.

For an object \(\pi\) in that missing notation satisfying \(r(\pi)>0\), write \(\mathcal D_{\pi}(s)\) for the formal Dirichlet series that the source says can be constructed via Hecke algebras and has Euler multiplication. This notation merely names the series mentioned in the source; its coefficients, Euler factors, and normalization are not supplied.

The source asks whether there exists an algebraic Hasse--Weil variety \(X\) whose zeta function is this series, schematically
\[
\exists X\qquad \zeta_X(s)=\mathcal D_{\pi}(s).
\]
This is an unresolved statement schema: the admissible class and base of \(X\), the convention for \(\zeta_X\), and the precise meaning of equality with the formal series are not specified. The record also does not determine the quantifier over \(\pi\). These omissions prevent a self-contained, fully quantified formulation. The source adds that there are several conjectures, without stating them.
\subsection{Short English statement}
For the objects \(\pi\) with \(r(\pi)>0\) referred to in Problem 8.82, the source asks whether a formal Dirichlet series constructed using Hecke algebras is the zeta function of an algebraic Hasse--Weil variety. The missing preceding notation, unspecified series, and unspecified variety and zeta-function conventions require clarification before this can be a precise standalone statement.
\subsection{Sources}
\begin{itemize}
\item Supplied local registry record: permanent UnsolvedMath ID \(700531\), exact source code \texttt{catalogue entry}, title \emph{Kourovka Notebook Problem 8.83}. The registry labels the problem \texttt{open}; this is supplied source metadata, not an independently verified mathematical or literature assessment.
\item \emph{The Kourovka Notebook (Unsolved Problems in Group Theory)}, Problem 8.83, proposed by H. Helling, 8th issue (1982). The supplied background identifies the 21st-issue edition (2026), \texttt{arXiv:1401.0300v48}. The registry describes the collection as edited by E. I. Khukhro and V. D. Mazurov. These bibliographic details have not been independently checked.
\item G. Shimura, \emph{Introduction to the arithmetic theory of automorphic functions}, Princeton Univ. Press, 1971. This is the reference supplied for the Hecke-algebra construction; no page, theorem, or explicit construction is identified.
\item Problem 8.82 is explicitly invoked for notation, but its statement is absent from the supplied record.
\end{itemize}
\end{document}
